\section{Preliminaries}
\label{sec:preliminaries}

\noindent\textbf{Visual Autoregressive Modeling.}
Visual AutoRegressive modeling (VAR)~\cite{VAR} redefines image generation as a
coarse-to-fine process: rather than predicting tokens one by one in raster-scan
order, VAR first generates a rough low-resolution representation of the image and
progressively refines it at finer scales. Formally, VAR consists of a multi-scale tokenizer and
an autoregressive transformer. The tokenizer encodes an image $I$ into a hierarchy
of $K$ discrete token maps $\{R_1, R_2, \ldots, R_K\}$ at progressively finer
resolutions, where $R_k \in \mathbb{Z}^{h_k \times w_k}$ and
$h_1 < h_2 < \cdots < h_K$. The transformer then autoregressively predicts the
token map at each subsequent scale conditioned on all coarser scales and text
embeddings $\Psi(t)$:
\begin{equation}
    p(R_1, \ldots, R_K) = \prod_{k=1}^{K} p\!\left(R_k \mid R_1, \ldots,
    R_{k-1},\, \Psi(t)\right)
\end{equation}
All tokens within a scale are predicted in parallel, and each token in $R_k$
corresponds to a spatially well-defined patch at resolution $(h_k, w_k)$. This
spatial locality is the structural property that makes token-level selective
editing geometrically meaningful in our framework.

\noindent\textbf{Infinity and Bitwise Tokenization.}
Infinity~\cite{Infinity} extends VAR by replacing the conventional index-wise
tokenizer with Binary Spherical Quantization (BSQ)~\cite{BSQVAE}, which projects
each continuous residual vector $z_k \in \mathbb{R}^d$ onto the unit hypersphere
and binarizes each dimension:
\begin{equation}
    q_k = \frac{1}{\sqrt{d}}\,\operatorname{sign}\!\left(\frac{z_k}{\|z_k\|}
    \right)
\end{equation}
This formulation scales the vocabulary to effectively infinite size while keeping
computational cost linear in $d$. We adopt Infinity-2B as the backbone of our
framework. Critically for our method, BSQ encoding is \textit{self-consistent}: encoding
a real image into its bitwise token representation and decoding it back introduces
no additional distortion beyond the tokenizer's inherent reconstruction error, meaning that transplanting source tokens into unedited regions
guarantees pixel-faithful background preservation by construction.

\noindent\textbf{Two source-information channels.}
Throughout this paper we distinguish two ways in which the source image enters the
generation process, because they behave very differently. The \textit{continuous}
channel injects the VAE encoder output $\mathcal{E}(I_s)$ (before quantization)
into the accumulated residual maintained across scales, and is used only to
reproduce the source image faithfully while its cross-attention is recorded. The
\textit{discrete} channel supplies the per-scale bitwise tokens
$\{r_k^s\} = \mathcal{Q}_{\mathrm{BSQ}}(\mathcal{E}(I_s))$, and is the only channel
that participates in token substitution. Every token written back into the output
therefore comes from the BSQ quantization of the source image itself, never from an
intermediate generation --- this is what makes background preservation exact rather
than approximate.
